\section{Experimental protocol}\label{sec:protocol}
\subsection{Comparators and adaptation scope}
Rolling calibration uses causally received absolute normalised residuals by sensor and horizon. The ACI adaptation maintains a horizon-specific significance controller. At a feedback arrival bin it updates once using the spatial batch's mean interval miss; there is no update before release and none for a permanently lost outcome. This batching and selective outcome suppression differ from the source protocol, so the source coverage guarantee is not inherited \citep{hallberg2024}.

CoRel is evaluated as a documented project adaptation on the FAR-GW prediction caches. Its architecture and relational calibration follow the official implementation, but the partitioning, multihorizon targets and reporting replay follow this study. These results compare calibration on identical frozen forecasts, not the scores reported on the source paper's splits \citep{cini2025}. CoRel is not extrapolated to the other backbones.

FAR-GW adapts Graph WaveNet with the seven causal channels and noncrossing quantile head. STAEformer is an independent paper-based implementation because the audited official repository did not provide a licence file. STID and DCRNN are documented adaptations of their licensed references. Each backbone is trained with seeds 11, 22 and 33; fault seeds are 101, 202 and 303. Table~\ref{tab:settings} records shared settings and calibrated-method parameters. Per-backbone configurations and source audits accompany the release.

\subsection{Reporting conditions, controls and stress tests}
Table~\ref{tab:conditions} defines the six reporting conditions. C3 combines correlated input outages with delayed feedback and backlog release. C4 imposes permanent feedback loss within the affected region. The baseline connected outage affects 10\% of sensors for six bins, with two events per day and a three-bin backlog release. C4 loses half of the within-outage feedback. Physical target validity remains unchanged.

The matched-rate independent control randomises affected sensor identities while retaining severity. It asks whether a method succeeds under the independent process; it does not by itself estimate a causal correlated-minus-independent treatment effect. The one-factor stress grid changes duration, connected fraction, additional backlog delay or within-outage loss while holding other settings fixed. Two further tests impair initial calibration and introduce additional low-speed feedback loss. The latter uses a training-derived 20th-percentile speed threshold and suppresses qualifying outcomes with additional probability 0.5. Truth enters only the evaluator-side fault generator. This is missing-not-at-random (MNAR) sensitivity, not an identifiable correction procedure.

\subsection{Endpoints and paired inference}
For interval $[L,U]$, outcome $y$ and miscoverage $\alpha$, the interval score is
\begin{equation}
\mathrm{IS}_{\alpha}(L,U;y)=U-L+\frac{2}{\alpha}(L-y)\mathbf1(y<L)+\frac{2}{\alpha}(y-U)\mathbf1(y>U).
\end{equation}
Lower scores reward narrow intervals while penalising missed targets. Coverage is the fraction of originally valid targets inside the interval; width is $U-L$. The primary endpoint averages the 90\% score equally over 15-, 30- and 60-minute horizons. Paired differences are comparator minus \FAR{}, so positive effects favour \FAR{}. The registered rule requires lower 95\% bounds above zero against both rolling and ACI, together with aggregate coverage at least 0.88. Nominal coverage remains 0.90; the 0.88 floor is an evaluation tolerance, not a revised nominal target.

SUMO inference averages model and fault seeds within each episode, then resamples complete episodes 2,000 times. Real-data inference averages model seeds within each fault seed and uses one-day moving blocks, stratified by fault seed. Model-seed repetitions are not independent temporal replications. Real-data audits use systematically selected hourly origins and affected sensors through 12 post-event bins; the buffer cap controls cost. Thus the real-data results do not represent every five-minute network-wide forecast. Stress summaries average origin-level scores, whereas descriptive runtime quality uses pooled valid targets; these estimands are not substituted for one another.

Seven component ablations retain the same point and scale caches: no spatial pooling, no context matching, no target-age weighting, no stale blend, no inflation, arrival-time freshness and frozen calibration only. The arrival-time diagnostic changes age weighting and freshness gap but retains target-based buffer eligibility and timestamp-block support. Frozen-only uses the same archive and fallback without live adaptation or inflation. Full-reference runs are verified and reused. Ablation intervals are exploratory and unadjusted for 42 comparisons; their signs are descriptive evidence rather than family-wise confirmatory tests.

